% GENERATED FILE. Edit canonical lessons, then run npm run generate:books.
% canonical-source-hash: fnv1a64:14d136191154e92e
% canonical-lessons: FR-C307-degat, FR-C307-endommager, FR-C307-assurance, FR-C307-sinistre, FR-C307-rayure, FR-R307-revision-300-307

\chapter{Les dégâts}
\label{ch:fr-les-degats}
\hlchaptermodality{307}

\begin{chapteropening}
\textbf{By the end of this chapter:} \emph{I can talk about damage after an accident: the damage, damaging, insurance, the people hit by a disaster and a scratch.}

Complete the last lesson of chapter 307: I can talk about damage after an accident: the damage, damaging, insurance, the people hit by a disaster and a scratch.
\end{chapteropening}

\section[le dégât]{le dégât — the damage}

\label{lesson:FR-C307-degat}



\begin{quote}
Before the new one: say the French for to bite, then the French for the bite.
\end{quote}

\subsection*{You'll want to know: le dégât}
\textbf{le dégât} — \textquotedblleft{}the damage\textquotedblright{}.

\begin{quote}
L'orage a fait beaucoup de dégâts.
\end{quote}

\begin{quote}
\emph{The storm did a lot of damage.}
\end{quote}

Usually plural: \emph{les dégâts}.

\begin{cousinweb}
\emph{dégât} is built on Old French \emph{degaster} \textquotedblleft{}to lay waste\textquotedblright{}, usually traced to Latin \emph{dēvastāre} (EtymDB lists no parent), built on \emph{vastus} \textquotedblleft{}empty, waste\textquotedblright{} (Lewis \& Short).
\end{cousinweb}

\begin{grammarlens}[title={What you've built}]
One more word for telling, explaining and describing at B1.
\end{grammarlens}

\subsection*{Your turn}
\begin{itemize}
\raggedright
  \item \emph{Say it:} \emph{le dégât}
  \item \emph{Say it:} \emph{le dégât}, once more
  \item \emph{Say it:} a sentence of your own with \emph{le dégât}
  \item \emph{Recall:} say the French for to bite, then the French for the bite, then say \emph{le dégât} again
\end{itemize}

\begin{tcolorbox}[breakable,colback=teal!4,colframe=teal!35!black,title={Before you move on}]
What does le dégât mean? (\textquotedblleft{}The damage\textquotedblright{}.) Say it once more.
\end{tcolorbox}
\section[endommager]{endommager — to damage}

\label{lesson:FR-C307-endommager}



\begin{quote}
Before the new one: say the French for the bite, then the French for the damage.
\end{quote}

\subsection*{You'll want to know: endommager}
\textbf{endommager} — \textquotedblleft{}to damage\textquotedblright{}.

\begin{quote}
L'eau a endommagé la cuisine.
\end{quote}

\begin{quote}
\emph{The water damaged the kitchen.}
\end{quote}

For things, not people.

\begin{cousinweb}
\emph{endommager} is built on \emph{dommage} \textquotedblleft{}harm\textquotedblright{}, from Old French \emph{damage}, which goes back to Latin \emph{damnum} \textquotedblleft{}loss\textquotedblright{} (Lewis \& Short).
\end{cousinweb}

\begin{grammarlens}[title={What you've built}]
One more word for telling, explaining and describing at B1.
\end{grammarlens}

\subsection*{Your turn}
\begin{itemize}
\raggedright
  \item \emph{Say it:} \emph{endommager}
  \item \emph{Say it:} \emph{endommager}, once more
  \item \emph{Say it:} a sentence of your own with \emph{endommager}
  \item \emph{Recall:} say the French for the bite, then the French for the damage, then say \emph{endommager} again
\end{itemize}

\begin{tcolorbox}[breakable,colback=teal!4,colframe=teal!35!black,title={Before you move on}]
What does endommager mean? (\textquotedblleft{}To damage\textquotedblright{}.) Say it once more.
\end{tcolorbox}
\section[l'assurance]{l'assurance — the insurance; self-confidence}

\label{lesson:FR-C307-assurance}



\begin{quote}
Before the new one: say the French for the damage, then the French for to damage.
\end{quote}

\subsection*{You'll want to know: l'assurance}
\textbf{l'assurance} — \textquotedblleft{}the insurance; self-confidence\textquotedblright{}.

\begin{quote}
L'assurance va tout payer.
\end{quote}

\begin{quote}
\emph{The insurance is going to pay for everything.}
\end{quote}

\emph{Avoir de l'assurance} — to be self-assured.

\begin{cousinweb}
\emph{assurance} is built on \emph{assurer} \textquotedblleft{}to make safe\textquotedblright{}, which is inherited from Vulgar Latin \emph{assēcūrāre}, built on Latin \emph{sēcūrus} \textquotedblleft{}safe, free from care\textquotedblright{} (Lewis \& Short).
\end{cousinweb}

\begin{grammarlens}[title={What you've built}]
One more word for telling, explaining and describing at B1.
\end{grammarlens}

\subsection*{Your turn}
\begin{itemize}
\raggedright
  \item \emph{Say it:} \emph{l'assurance}
  \item \emph{Say it:} \emph{l'assurance}, once more
  \item \emph{Say it:} a sentence of your own with \emph{l'assurance}
  \item \emph{Recall:} say the French for the damage, then the French for to damage, then say \emph{l'assurance} again
\end{itemize}

\begin{tcolorbox}[breakable,colback=teal!4,colframe=teal!35!black,title={Before you move on}]
What does l'assurance mean? (\textquotedblleft{}The insurance; self-confidence\textquotedblright{}.) Say it once more.
\end{tcolorbox}
\section[le sinistré]{le sinistré — the disaster victim}

\label{lesson:FR-C307-sinistre}



\begin{quote}
Before the new one: say the French for to damage, then the French for the insurance.
\end{quote}

\subsection*{You'll want to know: le sinistré}
\textbf{le sinistré} — \textquotedblleft{}the disaster victim\textquotedblright{}.

\begin{quote}
Les sinistrés ont dormi dans une école.
\end{quote}

\begin{quote}
\emph{The people hit by the disaster slept in a school.}
\end{quote}

Also an adjective: \emph{une région sinistrée}.

\begin{cousinweb}
\emph{sinistré} is built on \emph{sinistre} \textquotedblleft{}a disaster\textquotedblright{}, from Latin \emph{sinister} \textquotedblleft{}left, unlucky\textquotedblright{} (Lewis \& Short).
\end{cousinweb}

\begin{grammarlens}[title={What you've built}]
One more word for telling, explaining and describing at B1.
\end{grammarlens}

\subsection*{Your turn}
\begin{itemize}
\raggedright
  \item \emph{Say it:} \emph{le sinistré}
  \item \emph{Say it:} \emph{le sinistré}, once more
  \item \emph{Say it:} a sentence of your own with \emph{le sinistré}
  \item \emph{Recall:} say the French for to damage, then the French for the insurance, then say \emph{le sinistré} again
\end{itemize}

\begin{tcolorbox}[breakable,colback=teal!4,colframe=teal!35!black,title={Before you move on}]
What does le sinistré mean? (\textquotedblleft{}The disaster victim\textquotedblright{}.) Say it once more.
\end{tcolorbox}
\section[la rayure]{la rayure — the scratch; a stripe}

\label{lesson:FR-C307-rayure}



\begin{quote}
Before the new one: say the French for the insurance, then the French for the disaster victim.
\end{quote}

\subsection*{You'll want to know: la rayure}
\textbf{la rayure} — \textquotedblleft{}the scratch; a stripe\textquotedblright{}.

\begin{quote}
Il y a une longue rayure sur la voiture.
\end{quote}

\begin{quote}
\emph{There is a long scratch on the car.}
\end{quote}

Also a stripe: \emph{une chemise à rayures}.

\begin{cousinweb}
\emph{rayure} is built on \emph{rayer} \textquotedblleft{}to scratch, to stripe\textquotedblright{}, from \emph{raie} \textquotedblleft{}a line\textquotedblright{}, which EtymDB traces through Old French \emph{roie} to Gaulish \emph{rica} \textquotedblleft{}a furrow\textquotedblright{}; the details are not fully settled (EtymDB lists no parent).
\end{cousinweb}

\begin{grammarlens}[title={What you've built}]
One more word for telling, explaining and describing at B1.
\end{grammarlens}

\subsection*{Your turn}
\begin{itemize}
\raggedright
  \item \emph{Say it:} \emph{la rayure}
  \item \emph{Say it:} \emph{la rayure}, once more
  \item \emph{Say it:} a sentence of your own with \emph{la rayure}
  \item \emph{Recall:} say the French for the insurance, then the French for the disaster victim, then say \emph{la rayure} again
\end{itemize}

\begin{tcolorbox}[breakable,colback=teal!4,colframe=teal!35!black,title={Before you move on}]
What does la rayure mean? (\textquotedblleft{}The scratch; a stripe\textquotedblright{}.) Say it once more.
\end{tcolorbox}
\section[(dialogue)]{Retour sur les chapitres 300–307}

\label{lesson:FR-R307-revision-300-307}



\begin{quote}
Say the French for the emergency services and to rescue.
\end{quote}

\subsection*{Your turn}
Say these aloud:

\begin{itemize}
\raggedright
  \item the emergency services, to rescue, to resuscitate, the tourniquet, the stretcher
  \item the alert, to alert, the defibrillator, the emergency exit, to evacuate
  \item the fire, the fire extinguisher, the explosion, to explode, the spark
  \item the earthquake, the avalanche, the rockfall, to collapse, the rubble
  \item drowning, to drown, the rescuer, the buoy, the lifeguard
  \item the shipwreck, to crash, the survivor, to survive, the catastrophe
  \item to choke, the poisoning, to sting, to bite, the bite
  \item the damage, to damage, the insurance, the disaster victim, the scratch
\end{itemize}

\begin{tcolorbox}[breakable,colback=teal!4,colframe=teal!35!black,title={Before you move on}]
The emergency services? (\textbf{les secours}.) To rescue? (\textbf{secourir}.) To resuscitate? (\textbf{réanimer}.) The tourniquet? (\textbf{le garrot}.) The stretcher? (\textbf{la civière}.) The alert? (\textbf{l'alerte}.) To alert? (\textbf{alerter}.) The defibrillator? (\textbf{le défibrillateur}.) The emergency exit? (\textbf{la sortie de secours}.) To evacuate? (\textbf{évacuer}.) The fire? (\textbf{l'incendie}.) The fire extinguisher? (\textbf{l'extincteur}.) The explosion? (\textbf{l'explosion}.) To explode? (\textbf{exploser}.) The spark? (\textbf{l'étincelle}.) The earthquake? (\textbf{le tremblement de terre}.) The avalanche? (\textbf{l'avalanche}.) The rockfall? (\textbf{l'éboulement}.) To collapse? (\textbf{effondrer}.) The rubble? (\textbf{les décombres}.) Drowning? (\textbf{la noyade}.) To drown? (\textbf{noyer}.) The rescuer? (\textbf{le sauveteur}.) The buoy? (\textbf{la bouée}.) The lifeguard? (\textbf{le maître-nageur}.) The shipwreck? (\textbf{le naufrage}.) To crash? (\textbf{écraser}.) The survivor? (\textbf{le rescapé}.) To survive? (\textbf{survivre}.) The catastrophe? (\textbf{la catastrophe}.) To choke? (\textbf{étouffer}.) The poisoning? (\textbf{l'intoxication}.) To sting? (\textbf{piquer}.) To bite? (\textbf{mordre}.) The bite? (\textbf{la morsure}.) The damage? (\textbf{le dégât}.) To damage? (\textbf{endommager}.) The insurance? (\textbf{l'assurance}.) The disaster victim? (\textbf{le sinistré}.) The scratch? (\textbf{la rayure}.)
\end{tcolorbox}
